\section{References}
\label{sec:rational-bib}

\paragraph*{Rational functions and relations.}
To the best of my knowledge, the rational relations first appeared  in~\cite[Definition 15]{RabinScott59}  under the name of ``two-tape one-way automata''. The deterministic case of these automata does not correspond to the rational functions, since it describes a model where a deterministic automaton has access to both the input and output strings. To the best of my knowledge, the functional fragment first appears in \cite[Section II]{schutzenberger1961} under the name of ``transductions'', where the corresponding model is the same as what we call a bimachine, see~\cref{def:bimachine}.   Another early reference is 
\cite{elgotMezei1965}, which uses introduces the approach that is followed by this book: we first introduce rational relations, and then isolate the functions as a special case. In~\cite{elgotMezei1965}, the relations and functions are called ``recognised by generalised finite automata''. 


\paragraph*{Rational and recognisable subsets of a monoid.} None of the references discuss above use the name ``rational'', so we should explain why we use this name. It stems from a more general terminology, about rational and recognisable subsets of a monoid, which is popular in the French algebraic school. Let us explain this terminology. 

In the automata literature, the name rational seems to appear for the first time in Sch\"utzenberger's paper about (what is now known as) weighted automata~\cite[Section III.A]{Schutzenberger1961OnTheDefinition}, where it is applied to describe (what was previously known as) regular expressions on (non-commutative) power series. The name rational is particularly well motivated in the setting of power series, which is easiest to see for the monoid   $\field(x)$, which consists of fractions of polynomials over a field $\field$ that have one variable $x$, modulo common factors. This monoid is the least subset of all series (i.e.~functions from monomials $x^n$ to field coefficients) which contains all polynomials, is closed under addition and multiplication, and the following version of Kleene star:
        \begin{align*}
        1 + x + x^2 + \cdots + x^n + \cdots = \frac{1}{1-x}.
        \end{align*}
In a tradition dating back to at least Euler,  fractions of polynomials as in $\field(x)$ are called rational functions, which explains why this submonoid of the monoid of all series is called the rational series, and more generally why regular expressions are called rational. This is the terminology typically used in the French algebraic tradition. One early example which uses this name for a monoid which is not the free monoid is~\cite{eilenbergSchutzenberger1969} which talks about rational subsets of a commutative monoid -- for example in the commutative monoid $\Nat^d$ with component-wise additions, the rational subsets are exactly the semilinear sets. 

Let us now move to the recognisable subsets of a monoid, which are those that are recognised by a homomorphism into a finite monoid (this is the same as the congruence description from \cref{def:rational-recognisable-subsets}). The idea to use monoid to recognise languages goes back to Sch\"utzenberger, who defined the  syntactic monoid of a language in~\cite[p.~10]{schutzenberger1956}. An older reference which predates this is \cite[Theorem 3]{dubreil1941}, which introduces the syntactic congruence of a subset of a monoid (not necessarily a free monoid, and the congruence is a right congruence, as in a \dfa). Under the above terminology, the Kleene Theorem states that for the free monoid $A^*$, the rational subsets are exactly the recognisable subsets; for example this is the statement used \cite[Chapter III, Theorem 2.1]{berstel1979}. A more detailed discussion of this topic, with ample references, can be found in \cite{Sakarovitch2021Automata}. 

Let us now return to the use of the name rational for string-to-string functions and relations. As we have observed in the beginning of this section, none of the early papers use the name rational for string-to-string functions and relations; this includes Sch\"utzenberger himself. As far as I can tell, this terminology starts to appear around the time of Eilenberg's book~\cite[Chapter VII]{eilenberg1974}, maybe it is due to Eilenberg; Sch\"utzenberger seems to say so~\cite[Introduction]{reutenauer1991minimization}. Also Eilenberg introduced the name bimachine, and proved that it is equivalent to the rational functions, see~\cite[Chapter XI, Theorem 7.1]{eilenberg1974}.


\paragraph*{The equivalence problem.} Undecidability of the equivalence problem for rational relations was first proved in
\cite[p.410]{Griffiths1968}, although a related result on the undecidability of emptiness for intersection of rational relations can already be found in~\cite[Theorem 18]{RabinScott59}. The latter result  is our \cref{exer:rational-relations-intersection-undecidable}. 

The equivalence problem for rational functions was shown decidable in~\cite[Theorem 1]{blattner1977single}. This is deduced from a more general result, namely functionality for rational relations (two rational functions are equal if and only if their union is a functional rational relation). The functionality problem for relations (and hence also the equivalence problem for functions) is solved in~\cite{blattner1977single} using a a pumping arugment.    
 According to the editors comments in Sch\"utzenberger's collected works~\cite[p.4]{schutzenberger_oeuvres_tome08}, the functionality problem for relations is also implicitly solved in~\cite{schutzenberger1976}. 
 
 None of these decidability arguments above use a reduction to weighted automata, as we do here. 
    This kind of reduction traces back to \cite[Theorem 1]{AlbertLawrence1985} which proves a conjecture of Ehrenfeucht about string-to-string homomorphisms by embedding strings into the metabelian group and then applying the Hilbert Basis Theorem. The embedding itself is credited to Mal'cev~\cite{Malcev1953Nilpotent}. Of relevance to this book is also \cite[Theorem 1]{AlbertLawrence1985}, which uses the same approach to decide equivalence for HD0L equivalences; the latter equivalence is the same as copyful \sst's from \cref{exer:sst-copyful-sst} and it is the most powerful decidability result about equivalence that appears in this booi. A later result that is closer to our book because it uses weighted automata is  \cite{HarjuKarhumaki1991}. It is still not exactly the same as our book, because (a) the source of the reduction is different (instead of equivalence for rational string-to-string functions, they studied equivalence for  (tuple-of-string)-to-Boolean functions; and (b) the target of the reduction is different (instead of weighted automata over a field, they used weighted automata over a division ring). Also, since \cite{HarjuKarhumaki1991}  uses weighted automata over division ring, it  needs to extend Sch\"utzenberger's equivalence argument from a field to a division ring. The simple reduction from this book, which uses rational numbers to encode strings, appears in~\cite[p. 21:4]{Seidl2018} and is also discussed in~\cite[p. 6]{Bojanczyk2019Hilbert}.

\paragraph*{Machine independent characterisations.} As remarked in the beginning of this book, the characterisation of Mealy machines in \cref{thm:mealy-machine-independent} is in fact essentially the same as the original formulation of the Myhill-Nerode Theorem~\cite[Lemma]{nerode1958}. The characterisation of sequential functions in ~\cref{thm:sequential-function-independent} is from~\cite[p. 381]{GinsburgRose1966}. The characterisation of subsequential  functions in \cref{thm:subsequential-functions} is from~\cite[Proposition 3.4]{choffrut1977}. Finally, the characterisation of rational functions  from \cref{sec:rational-myhill-nerode} is from~\cite[Theorem 1 on p. 674]{reutenauer1991minimization}.